
%(BEGIN_QUESTION)
% Copyright 2006, Tony R. Kuphaldt, released under the Creative Commons Attribution License (v 1.0)
% This means you may do almost anything with this work of mine, so long as you give me proper credit

Anta at vi må måle konsentrasjonen av stoff {\it X} i en prosessgass-strøm. Hvis vi vet at {\it X} danner ioner ved forbrenning, betyr dette at vi bare kan koble utgangen fra en flammeionisasjonsdetektor (FID) til en indikator og få garantert måling av konsentrasjonen til stoff {\it X}? Hvorfor eller hvorfor ikke?

\underbar{file i00927}
%(END_QUESTION)





%(BEGIN_ANSWER)

Nei, fordi vi ikke vet hvilke {\it andre} stoffer i prosessstrømmen som også danner ioner ved forbrenning, og målingen trenger derfor ikke å representere konsentrasjonen av {\it X} alene. Vi kan likevel bruke FID til å måle dette stoffet dersom det er detektor på enden av en {\it kromatograf}, der stoff {\it X} forventes å komme ut av kolonnen på et bestemt tidspunkt. Da gjøres FID-ens ikke-spesifikke deteksjon spesifikk gjennom den selektive tidsforsinkelsen i kromatografkolonnen.

%(END_ANSWER)





%(BEGIN_NOTES)


%INDEX% Measurement, analytical: applications of different technologies

%(END_NOTES)

